\documentclass[11pt,a4paper]{article}
\usepackage[utf8]{inputenc}
\usepackage[spanish]{babel}
\usepackage[margin=2.5cm]{geometry}
\usepackage{xcolor}
\definecolor{navy}{RGB}{31,78,121}
\begin{document}
\begin{center}
{\Large\bfseries\color{navy} ANEXO F: PROPUESTA DE CONTROL GEODÉSICO Y TOPOGRÁFICO}\\[0.5em]
{\large MCH Río Blanco — 100 kW — Prefactibilidad}\\[0.5em]
{\normalsize Puno, Perú — 16 de octubre de 2026}
\end{center}
\section*{Estado del levantamiento}
Este anexo no acredita trabajos GNSS, monumentación, cotas, precisión ni certificación del Instituto Geográfico Nacional. No se adjuntan observaciones crudas, libretas, actas, certificados de estación base, informes de ajuste ni fichas oficiales. Los puntos y coordenadas incluidos en otros archivos son referencias digitales provisionales y no deben usarse para replanteo.
\section*{Trabajo requerido para la siguiente etapa}
Un profesional competente deberá: (1) confirmar el marco horizontal y vertical aplicable y su transformación; (2) definir y documentar vértices de control y su vínculo a estaciones oficiales; (3) levantar captación, conducción, cámara de carga, tubería, casa de máquinas, accesos y cruce de cauce; (4) registrar datos crudos, equipos, calibraciones, tiempos de observación, alturas de antena y condiciones; (5) ajustar la red y reportar cierres, incertidumbre y control independiente; (6) contrastar cotas y distancias con nivelación y observaciones redundantes; y (7) entregar planos, fichas, archivos digitales y metadatos con revisión profesional.
\section*{Parámetros todavía por validar}
La zona UTM, datum, época, modelo geoidal, conversión a elevación ortométrica, escala y tolerancias deben definirse de acuerdo con el levantamiento y la cartografía oficial. No se adoptan aquí ondulación geoidal, factor combinado ni orden geodésico como resultados medidos.
\end{document}
